% ==============================================================================
% Plantilla Institucional de Trabajo Terminal — UPIIZ IPN
% Archivo: capitulos/08-validacion.tex
% ==============================================================================

\chapter{\tituloValidacion}
\label{cap:validacion}

En este capítulo se describen los protocolos experimentales planificados, las métricas cuantitativas que se utilizarán y los criterios de aceptación para evaluar el desempeño integral del robot rover en el entorno de pruebas de terreno irregular.

% ------------------------------------------------------------------------------
\section{Métricas cuantitativas de evaluación de la misión}
\label{sec:validacion-metricas-rover}

Para evaluar rigurosamente el desempeño de cada subsistema y la eficacia global de la misión de exploración, se emplearán las siguientes métricas:

\begin{enumerate}
    \item \textbf{Error de posición euclidiano en navegación sin GPS:}
    \begin{equation}
        e_p(t) = \sqrt{\big(x_r(t) - x(t)\big)^2 + \big(y_r(t) - y(t)\big)^2}
        \label{eq:error-posicion-rover}
    \end{equation}
    \item \textbf{Tasa de éxito en la detección de muestras geológicas ($\eta_{\text{det}}$):}
    \begin{equation}
        \eta_{\text{det}} = \frac{N_{\text{muestras detectadas}}}{N_{\text{muestras presentes en el campo}}} \times 100\,\%
        \label{eq:tasa-deteccion}
    \end{equation}
    \item \textbf{Tasa de éxito del mecanismo de recolección ($\eta_{\text{rec}}$):}
    \begin{equation}
        \eta_{\text{rec}} = \frac{N_{\text{muestras recolectadas y aseguradas}}}{N_{\text{intentos de recolección}}} \times 100\,\%
        \label{eq:tasa-recoleccion}
    \end{equation}
    \item \textbf{Eficiencia global de cumplimiento de misión ($\eta_{\text{misión}}$):}
    \begin{equation}
        \eta_{\text{misión}} = \frac{N_{\text{misiones completadas exitosamente}}}{N_{\text{misiones ejecutadas}}} \times 100\,\%
        \label{eq:eficiencia-mision}
    \end{equation}
\end{enumerate}

% ------------------------------------------------------------------------------
\section{Protocolo y criterios de validación experimental en pista}
\label{sec:validacion-protocolo-rover}

En la \tabref{tab:comparativa-validacion} se resumen los criterios de aceptación cuantitativos que se verificarán durante los ensayos de campo en el escenario experimental de terreno irregular.

\begin{table}[htbp]
    \centering
    \caption{Criterios de validación y especificaciones de desempeño por verificar en el rover.}
    \label{tab:comparativa-validacion}
    \begin{tabularx}{\textwidth}{p{3.8cm} c p{3.8cm} X}
        \toprule
        \textbf{Prueba / Subsistema} & \textbf{Especificación} & \textbf{Método de Medición} & \textbf{Criterio de Aceptación} \\
        \midrule
        Movilidad en pendiente & $\ge \ang{15}$ de inclinación & Ensayos en rampa de tierra & Ascenso sin deslizamiento crítico \\
        Localización sin GPS & $e_p \le 5\,\%$ de trayecto & Contraste con marcas terrestres & Deriva acotada en circuito \\
        Detección visual de muestras & $\eta_{\text{det}} \ge 85\,\%$ & Registro de cámara y HMI & Reconocimiento a $\ge \qty{1.0}{\meter}$ \\
        Recolección y depósito & $\eta_{\text{rec}} \ge 80\,\%$ & Ensayos de manipulación & Captura, traslado y descarga \\
        Misión autónoma integral & $\eta_{\text{misión}} \ge 75\,\%$ & Ejecución completa sin GPS & Retorno exitoso a zona base \\
        \bottomrule
    \end{tabularx}
\end{table}

El cumplimiento de estos criterios permitirá demostrar la viabilidad, robustez y autonomía de la plataforma robótica desarrollada para misiones de exploración terrestre y lunar análoga.
